
Table \ref{tab:asQspiPer05} shows the input and outputs of the Qspi BPI.
\begin{table}[H]
\caption{QSPI BPI I/O}
\label{tab:asQspiPer05}
\centering
\begin{tabularx}{\textwidth}{|l |l |l |X|}
  \hline
  Pin & Direction & Width & Explanation \\
  \hline
  \hline
  clk\_i & in & 1 & System clock \\
  \hline
  rst\_i & in & 1 & System reset. Active high. \\
  \hline
  \multicolumn{4}{|l|}{Interface to peripheral kernel:} \\
  \hline
  ctrl\_reg\_o(4:0) & out & 5 & To QSPI-kernel. Defines and starts the QSPI FSM. \\
  \hline
  cmd\_reg\_o(x:0) & out & x & To QSPI-kernel. Defines the SPI command opcode. \\
  \hline
  addr\_reg\_o(x:0) & out & x & To QSPI-kernel. Defines the target flash address. \\
  \hline
  len\_reg\_o(x:0) & out & x & To QSPI-kernel. Defines the transfer length in bytes. \\
  \hline
  dummy\_reg\_o(x:0) & out & x & To QSPI-kernel. Defines the needed dummy cycles. \\
  \hline
  status\_reg\_i(3:0) & in & 4 & From QSPI-kernel. Delivers BUSY, DONE, ERROR and TIMEOUT. \\
  \hline
  tx\_data\_o(63:0) & out & 64 & To QSPI-kernel (TX-FIFO). Data to be send. \\
  \hline
  rx\_data\_i(63:0) & in & 64 & From QSPI-kernel (RX-FIFO). Data received. \\
  \hline
  \multicolumn{4}{|l|}{Interface to the bus:} \\
  \hline
  wbd\_s\_addr\_i(3:0) & in & 4 &  From WB-Bus. Peripherals addresses. \\
  \hline
  wbd\_s\_dat\_i(63:0) & in & 64 &  From WB-Bus. Peripherals data input. \\
  \hline
  wbd\_s\_dat\_o(63:0) & out & 64 &  To WB-Bus. Peripherals data output. \\
  \hline
  wbd\_s\_we\_i & in & 1 &  From WB-Bus. Peripherals write enable. \\
  \hline
  wbd\_s\_sel\_i(7:0) & in & 8 &  From WB-Bus. Peripherals byte select. Not used here. \\
  \hline
  wbd\_s\_stb\_i & in & 1 &  From WB-Bus. Valid bus cycle.  \\
  \hline
  wbd\_s\_ack\_o & out & 1 &  To WB-Bus. Error free bus transaction.   \\
  \hline
  wbd\_s\_cyc\_i & in & 1 &  From WB-Bus. High for complete bus cycle.  \\
  \hline
  \multicolumn{4}{|l|}{Other signals:} \\
  \hline
  qspi\_irq\_o & out & 1 &  QSPI interrupt.  \\
  \hline
\end{tabularx}
\end{table}
